\documentclass[11pt,a4paper]{article}
\usepackage[T1]{fontenc}\usepackage[utf8]{inputenc}\usepackage[english]{babel}
\usepackage{lmodern}\usepackage[a4paper,margin=2.5cm]{geometry}
\usepackage{amsmath,amssymb}\usepackage{siunitx}\usepackage{booktabs}
\usepackage{empheq}\usepackage[numbers,sort&compress]{natbib}
\usepackage[colorlinks=true,linkcolor=black,citecolor=black,urlcolor=blue]{hyperref}
\emergencystretch=1.5em
\newcommand{\crest}{\ensuremath{C}}
\newcommand{\dd}{\mathrm d}
\newcommand{\Au}{\ensuremath{A_\mathrm{u}}}
\newcommand{\Pprof}{\ensuremath{\tilde P_\mathrm{prof}}}
\title{Companion to the appendix:\\ full derivations}
\author{}\date{}
\begin{document}\maketitle
\noindent This document carries out, step by step, the derivations that the
appendix quotes as results. The section numbers are the ones the appendix
refers to. Nothing here is new physics; it exists so that every result in the
appendix can be checked without reconstructing it.
\tableofcontents

\section{Notation and conventions}\label{sec:not}
Tones are indexed $n=0,\dots,N-1$ per axis, with equidistant frequencies
$f_n=f_\mathrm{off}+n\,d$ ($d$ = tone spacing, not to be confused with the beat
order, which only appears in Sec.~\ref{sec:sinc}). The \emph{centred} tone index
\begin{equation}
\hat n:=n-\frac{N-1}{2}
\qquad\Longrightarrow\qquad f_n-\bar f=d\,\hat n
\label{eq:nhat}
\end{equation}
is used throughout, because it makes the frequency offset and the phase share
one factor. The complex envelope of the RF signal is
\begin{equation}
s(t)=\operatorname{Re}\!\left[A(t)e^{i2\pi\bar ft}\right],
\qquad
A(t)=\sum_{n}a_n\,e^{i(2\pi d\hat nt+\varphi_n)} ,
\label{eq:env}
\end{equation}
with $a_n$ the RF \emph{voltage} amplitudes. The profile parameters $r_x,r_y$
are \emph{power} ratios, so the edge tones carry
$\varrho=\sqrt{r}$ in voltage. $|A|$ has period $T_\mathrm{env}=1/d$; for even
$N$ the index $\hat n$ is half-integer and $A(t+T_\mathrm{env})=-A(t)$, but
$|A|$ is periodic either way, and only $|A|$ enters below.

\section{The centring criterion}\label{sec:centring}

\subsection{What ``centred'' has to mean}
The atom sits at $\mathbf r=0$, smeared over $\sigma\approx\SI{108}{\nano\meter}$
\citep{tuchendler2008}. The damaging situation is not a beating pattern but a
pattern that sits \emph{off to one side} during the pulse: the pulse area is
then larger on one edge of the position distribution than on the other, the
rotation angle depends on where the atom happened to be, and no calibration of
the mean repairs it. The requirement is therefore
\begin{equation}
I(-\mathbf r,t)=I(\mathbf r,t)\qquad\text{for \emph{all} }t,
\label{eq:centred}
\end{equation}
which implies $\nabla I|_0=0$ at all times.

\subsection{The three symmetries of the profile}
With $u$ the (even) single-spot mode and $\bar n=N-1-n$ the mirror index,
\begin{align}
\mathbf c_{\bar n\bar m}&=-\mathbf c_{nm} &&\text{(positions: the grid is centred)},\label{eq:symgeo}\\
a_{\bar n}&=a_n,\quad b_{\bar m}=b_m &&\text{(amplitudes: }[\varrho,1,\dots,1,\varrho]),\label{eq:symamp}\\
f_{\bar n}&=F-f_n,\quad g_{\bar m}=G-g_m &&\text{(frequencies)} .\label{eq:symfreq}
\end{align}
Equation~\eqref{eq:symfreq} is immediate from the raster:
$f_n+f_{\bar n}=2f_\mathrm{off}+(N-1)d=:F$, independent of $n$.

\subsection{Derivation}
Start from
$E(\mathbf r,t)=\sum_{n,m}a_nb_m\,u(\mathbf r-\mathbf c_{nm})
e^{-i2\pi(f_n+g_m)t}e^{i(\varphi_n+\psi_m)}$ and evaluate it at $-\mathbf r$,
relabelling the sum by the mirror index (a permutation, so nothing is lost).

\paragraph{Step 1 (positions; $u$ even).}
$-\mathbf r-\mathbf c_{\bar n\bar m}=-\mathbf r+\mathbf c_{nm}
=-(\mathbf r-\mathbf c_{nm})$, and $u$ even gives
$u(-\mathbf r-\mathbf c_{\bar n\bar m})=u(\mathbf r-\mathbf c_{nm})$: the same
envelope factor as in $E(\mathbf r,t)$.

\paragraph{Step 2 (amplitudes).} $a_{\bar n}b_{\bar m}=a_nb_m$, unchanged.

\paragraph{Step 3 (frequencies) --- the decisive one.}
\begin{equation}
e^{-i2\pi(f_{\bar n}+g_{\bar m})t}
=e^{-i2\pi(F+G)t}\;
\underbrace{e^{+i2\pi(f_n+g_m)t}}_{\text{the \emph{conjugate} of the original}} .
\end{equation}
The frequency mirror does not reproduce the time factor, it \emph{conjugates}
it, up to one global factor that is the same for every spot.

\paragraph{Step 4 (collect).} Comparing
\begin{align}
E(-\mathbf r,t)&=e^{-i2\pi(F+G)t}\sum_{n,m}a_nb_m\,u(\mathbf r-\mathbf c_{nm})
e^{+i2\pi(f_n+g_m)t}\,e^{i(\varphi_{\bar n}+\psi_{\bar m})},\\
E^{*}(\mathbf r,t)&=\sum_{n,m}a_nb_m\,u(\mathbf r-\mathbf c_{nm})
e^{+i2\pi(f_n+g_m)t}\,e^{-i(\varphi_{n}+\psi_{m})},
\end{align}
the two sums run over the same basis functions with the same real coefficients
and differ only in a constant phase per term. They are proportional --- and
hence $I(-\mathbf r,t)=I(\mathbf r,t)$ --- if and only if
$e^{i(\varphi_{\bar n}+\psi_{\bar m})}=e^{i\Theta}e^{-i(\varphi_n+\psi_m)}$
with a \emph{single} $\Theta$. Separating the $n$- and $m$-dependence,
\begin{empheq}[box=\fbox]{equation}
\varphi_n+\varphi_{N-1-n}=c\quad\forall n,
\qquad
\psi_m+\psi_{M-1-m}=c'\quad\forall m .
\label{eq:centring}
\end{empheq}

\subsection{Conjugate, not mirrored --- and the self-paired middle tone}
\label{sec:selfpair}
Equation~\eqref{eq:centring} is often misread as $\varphi_{\bar n}=\varphi_n$.
That is a different condition and it does not centre the pattern: the mirror
conjugates, so the phases must \emph{sum} to a constant. Numerically,
mirror-symmetric sets leave centroid offsets of \SIrange{44}{84}{\nano\meter}
against $\sigma=\SI{108}{\nano\meter}$; sets obeying \eqref{eq:centring} give
$0.00$~nm. The two coincide only for real phasors, $\varphi_n\in\{0,\pi\}$,
which is where the confusion comes from.

A natural objection is that for odd $N$ the middle tone has no partner. It
has: the map $n\mapsto\bar n=N-1-n$ is an involution, and involutions have
pairs \emph{and fixed points} ($N=3$: $0\leftrightarrow2$, $1\mapsto1$;
$N=4$: none; $N=5$: $2\mapsto2$). A fixed point is not unpaired, it is paired
with itself, and \eqref{eq:centring} holds for every $n$, so
\begin{equation}
2\varphi_{(N-1)/2}=c
\qquad\Longrightarrow\qquad
\varphi_{(N-1)/2}=\frac c2\ \text{ or }\ \frac c2+\pi .
\label{eq:midphase}
\end{equation}
The objection implicitly assumes each pair may carry its own constant. It may
not: $c$ \emph{is} the single global $\Theta$, and one $e^{i\alpha(t)}$ has to
serve all spots at once. The middle tone sits at $\bar f$, so its spot lies at
the mirror centre and maps onto itself.

\paragraph{Numerical check.} With $\varphi_x=(0,\psi,0)$ --- so $c=0$ --- and
$\varphi_y=(0,98,16,114)^\circ$ held fixed (which satisfies its own criterion,
$0+114=98+16$), the largest intensity gradient at the atom over a full beat
period is

\begin{center}\small
\begin{tabular}{crrl}
\toprule
$\psi$ & $\max_t|\partial_xI|/I$ & $\max_t|\partial_yI|/I$ &\\
\midrule
$0^\circ$ & $4\cdot10^{-10}$ & $4\cdot10^{-10}$ & criterion satisfied\\
$10^\circ$ & $8.8\cdot10^{-4}$ & $2.6\cdot10^{-4}$ &\\
$30^\circ$ & $2.5\cdot10^{-3}$ & $7.4\cdot10^{-4}$ &\\
$90^\circ$ & $5.1\cdot10^{-3}$ & $1.5\cdot10^{-3}$ &\\
$180^\circ$ & $4\cdot10^{-10}$ & $4\cdot10^{-10}$ & criterion satisfied\\
\bottomrule
\end{tabular}
\end{center}

\noindent The gradient vanishes at exactly the two values predicted by
\eqref{eq:midphase} and nowhere else; with $c=60^\circ$ the zeros move to
$30^\circ$ and $210^\circ$. Note the third column: violating the criterion in
$x$ also tilts the $y$ gradient. With a separable (Gaussian) envelope it would
not ($3\cdot10^{-10}$ in every row); the real Airy envelope is radial, hence
not separable, and couples the axes. The criterion must hold on \emph{both}.

\subsection{Structure of the admissible family}
Writing $\varphi_n=c/2+\chi_n$ turns \eqref{eq:centring} into
\begin{equation}
\chi_{\bar n}=-\chi_n ,
\label{eq:antisym}
\end{equation}
a global phase plus an antisymmetric part. Antisymmetric vectors in
$\mathbb R^N$ span $\lfloor N/2\rfloor$ dimensions, so the criterion halves the
phase freedom, from $N-1$ to $\lfloor N/2\rfloor$. One of those is always a
ramp and therefore crest-neutral (Sec.~\ref{sec:ramp}), leaving
$\lfloor N/2\rfloor-1$ crest-relevant dimensions --- zero for $N=3$.

\section{The crest factor}\label{sec:crest}

\subsection{The denominator does not depend on the phases}\label{sec:rms}
With $\crest=\max_t|s|/\mathrm{rms}_t[s]$ and the two time scales separated
($\langle s^2\rangle=\tfrac12\langle|A|^2\rangle$, while the peak of $|s|$ is
the peak of $|A|$),
\begin{equation}
\crest=\sqrt2\,\frac{\max_t|A(t)|}{\mathrm{rms}_t|A(t)|} .
\label{eq:crestenv}
\end{equation}
Expanding $\langle|A|^2\rangle$ gives a double sum whose off-diagonal terms are
phasor averages,
\begin{equation}
\bigl\langle e^{i[2\pi\Delta f\,t+\Delta\varphi]}\bigr\rangle_T
=e^{i\Delta\varphi}\,\frac{e^{i2\pi\Delta f\,T}-1}{i\,2\pi\Delta f\,T},
\label{eq:phasoravg}
\end{equation}
which vanishes \emph{exactly} when $\Delta f\,T$ is a non-zero integer: a
phasor completing whole turns averages to the centre of the circle
\emph{regardless of where it starts}, and that is precisely why
$\Delta\varphi$ drops out. With $T=T_\mathrm{env}=1/d$ every
$\Delta f=f_n-f_m$ is an integer multiple of $d$, so
\begin{empheq}[box=\fbox]{equation}
\mathrm{rms}_t|A|=\sqrt{\textstyle\sum_na_n^2}
\qquad\text{for \emph{every} phase set.}
\label{eq:rmsphase}
\end{empheq}
If the window is chosen shorter, \eqref{eq:phasoravg} no longer vanishes and
the rms becomes phase dependent: for the $x$ axis, using $T=1/\Delta f$ of the
neighbour pair instead of $1/d$, the rms varies between $1.548$ and $2.111$ for
different phase sets, against $1.714643$ always with the correct window.

\subsection{Bounds}
$|A(t)|\le\sum_na_n$ at all times (triangle inequality), and $|A|$ reaches at
least its own rms somewhere, so
\begin{equation}
\sqrt2\ \le\ \crest\ \le\ \sqrt2\,\frac{\sum_na_n}{\sqrt{\sum_na_n^2}}
\ \xrightarrow{\;a_n\equiv1\;}\ \sqrt{2N} .
\label{eq:bounds}
\end{equation}
The upper bound is attained exactly when all phasors align at some instant.
Which sets approach the lower bound is classical: Schroeder's quadratic phases
reach $\crest\approx1.7$ for large $N$ \citep{schroeder1970}, Kitayoshi's
construction is comparable \citep{kitayoshi1985}, and the attainable minimum is
studied systematically in \citep{boyd1986}. \emph{None of them satisfies
Eq.~\eqref{eq:centring}}, which is the whole tension.

\subsection{A ramp is a time shift, and therefore crest-neutral}\label{sec:ramp}
A ramp is $\varphi_n=\varphi_\mathrm{s}+\delta n$. The start value is common to
all tones and factors out; subtracting the mean
$\bar\varphi=\varphi_\mathrm{s}+\delta(N-1)/2$ leaves
\begin{equation}
\chi_n:=\varphi_n-\bar\varphi=\delta\,\hat n .
\label{eq:rampchi}
\end{equation}
This is not a replacement but a factorisation: with
$\varphi_n=\bar\varphi+\chi_n$,
\begin{equation}
A(t)=e^{i\bar\varphi}\underbrace{\sum_na_ne^{i(2\pi d\hat nt+\chi_n)}}
_{=:\tilde A(t)},
\qquad |A|=|\tilde A|,
\qquad \crest(A)=\crest(\tilde A) .
\end{equation}
Since $\hat{\bar n}=-\hat n$, Eq.~\eqref{eq:rampchi} satisfies
\eqref{eq:antisym}: \textbf{every ramp obeys the centring criterion
automatically.}

\paragraph{Does a shift exist?} The question is not how to define $\tau$ but
whether one number $\tau$ --- the same for all times and all tones --- can
satisfy $\tilde A_\delta(t)=\tilde A_0(t+\tau)$ for all $t$. Both sides expand
in the same linearly independent functions $e^{i2\pi d\hat nt}$, so
coefficients may be compared:
\begin{equation}
a_ne^{i\delta\hat n}=a_ne^{i2\pi d\hat n\tau}
\iff
2\pi d\,\hat n\,\tau\equiv\delta\,\hat n \pmod{2\pi}\quad\text{for every }n .
\label{eq:tausystem}
\end{equation}
This is a \emph{system}: $N$ conditions, one unknown. Generically unsolvable
--- but here the same factor $\hat n$ stands on both sides, cancels for
$\hat n\neq0$ and gives $0=0$ for $\hat n=0$, so all $N$ conditions collapse
into one:
\begin{empheq}[box=\fbox]{equation}
\tau=\frac{\delta}{2\pi d}
\qquad\Longrightarrow\qquad
\tilde A_\delta(t)=\tilde A_0(t+\tau) .
\label{eq:rampshift}
\end{empheq}
\emph{The content is the solvability, not the value of $\tau$.}

\paragraph{Converse.} For an arbitrary mean-free set $\chi_n$ the same
comparison demands $\tau=\chi_n/(2\pi d\hat n)$, a number that generally
depends on $n$. A common $\tau$ exists iff $\chi_n/\hat n$ is the same for all
tones, i.e.\ iff $\chi_n\equiv\delta\hat n$ modulo $2\pi$. \textbf{A phase set
is a time shift of the envelope if and only if it is a ramp.}

\paragraph{Why the crest factor cannot notice.} Both $\max_t$ and
$\mathrm{rms}_t$ run over a \emph{full} period, and $t'=t+\tau$ maps a full
period onto a full period; the \emph{set} of values of $|\tilde A|$ is
unchanged, only its ordering. Hence $\crest(\delta)=\crest(0)$ exactly.
Numerically, over $400\,000$ samples and $\delta=0^\circ\dots359^\circ$:
$\max|\tilde A|=2.96977156$, $\mathrm{rms}=1.71464282$,
$\crest=2.44942630$ for $N=3$, and $4.15406592/2.07846097/2.82648385$ for
$N=4$ --- agreement to eleven digits.

\paragraph{What a ramp does change.} The shift is real against an
\emph{absolute} time origin, and the trigger $t_0$ is one. At fixed $t_0=0$ the
pulse area over \SI{1}{\micro\second} drops to $0.825$ at
$\delta=12.5^\circ$ and to $0.043$ at $90^\circ$ of its value at $\delta=0$.
``Redundant with $t_0$'' means it must be considered \emph{together with} the
trigger, not that it is harmless. Across two axes, only \emph{matched} ramps
($\delta_x/d_x=\delta_y/d_y$, i.e.\ $\tau_x=\tau_y$) shift the whole
two-dimensional pattern: checked over a full beat period and the entire grid,
matched ramps give
$\max|I_\mathrm{ramp}(\mathbf r,t)-I_0(\mathbf r,t-\tau)|/\max I=8\cdot10^{-14}$,
a ramp on $x$ alone gives $0.28$.

\subsection{\boldmath$N=3$: the criterion forces the worst crest factor}
\label{sec:n3}
With $\theta:=2\pi dt$ and voltage amplitudes $a_0=a_2=\varrho$, $a_1=1$,
\begin{equation}
A(t)=\varrho e^{i(\varphi_0-\theta)}+e^{i\varphi_1}+\varrho e^{i(\varphi_2+\theta)} .
\end{equation}
The criterion requires $\varphi_0+\varphi_2=c$ and $2\varphi_1=c$, so the outer
phases lie symmetrically about $c/2$:
$\varphi_0=\tfrac c2-\chi$, $\varphi_2=\tfrac c2+\chi$,
$\varphi_1=\tfrac c2$ or $\tfrac c2+\pi$, with $\chi$ the only remaining
parameter. Factoring out $e^{ic/2}$,
\begin{empheq}[box=\fbox]{equation}
A(t)e^{-ic/2}
=\varrho e^{-i(\chi+\theta)}\pm1+\varrho e^{+i(\chi+\theta)}
=2\varrho\cos(\theta+\chi)\ \pm\ 1 .
\label{eq:n3env}
\end{empheq}
The whole time dependence now sits in a \emph{real} cosine, which sweeps
$[-1,1]$ over one period no matter what $\chi$ is --- $\chi$ only sets
\emph{when}. Maximising over $u:=\cos(\theta+\chi)\in[-1,1]$,
\begin{align}
\text{branch }+:&\quad\max_u|2\varrho u+1|=2\varrho+1 &&\text{at }u=+1,\\
\text{branch }-:&\quad\max_u|2\varrho u-1|=|-2\varrho-1|=2\varrho+1 &&\text{at }u=-1,
\end{align}
so $\max_t|A|=2\varrho+1=\sum_na_n$ for both branches and every $\chi,c$ ---
which is the triangle-inequality bound \eqref{eq:bounds}, valid for \emph{any}
phase set. The admissible family therefore always attains it.

\medskip
\noindent\fbox{\parbox{0.96\textwidth}{\textbf{With three tones the centring
criterion fixes the crest factor at the largest value the amplitudes allow, for
every $c$, every $\chi$ and both branches.} The reason is visible in
\eqref{eq:n3env}: the only free parameter appears \emph{inside} the cosine. It
can move \emph{when} the phasors align, not \emph{whether} they do.}}

\medskip
\noindent With $\varrho=\sqrt{r_x}=\sqrt{0.97}=0.9848858$,
\begin{equation}
\sum_na_n=2\varrho+1=2.9697716,\qquad
\sqrt{\textstyle\sum_na_n^2}=\sqrt{2r_x+1}=\sqrt{2.94}=1.7146428 ,
\end{equation}
--- note that the square sum contains the \emph{power} ratio linearly, the
square root cancelling --- hence $\crest_x=\sqrt2\cdot2.9697716/1.7146428
=2.4494263$. For exactly equal amplitudes this would be $\sqrt{2N}=\sqrt6
=2.4494897$; the working point lies $6\cdot10^{-5}$ below, so the two agree to
four digits without being the same number.

\paragraph{The family is the ramp family.} The two branches
$[0,c/2,c]$ and $[0,c/2+\pi,c]$ look distinct but are the two halves of one
circle: the ramp $\varphi_n=\delta n$ gives $c=2\delta$ and
$\varphi_1=\delta$, and as $\delta$ runs once through $[0,2\pi)$,
$c=2\delta\bmod2\pi$ runs through it \emph{twice} --- for $\delta<\pi$ with
$\varphi_1=c/2$, for $\delta\ge\pi$ with $\varphi_1=c/2+\pi$. Checked at
$\delta=0^\circ\dots315^\circ$ in $45^\circ$ steps: $\crest_x=2.4494$ in all
eight cases.

\subsection{\boldmath$N=4$, and the price of abandoning the criterion}
For four tones \eqref{eq:centring} leaves $\varphi_y=[0,b,c-b,c]$ with two
parameters; the ramp is $b=c/3$, and any other $b$ genuinely reshapes the
envelope. At $c=114^\circ$:

\begin{center}\small
\begin{tabular}{cccl}
\toprule
$b$ & $\varphi_y$ [deg] & $\max|A|$ & $\crest_y$\\
\midrule
$38^\circ$ & $[0,38,76,114]$ & $4.154$ & $2.8265$ (ramp, $=\sum a_n$)\\
$60^\circ$ & $[0,60,54,114]$ & $4.022$ & $2.7364$\\
$\mathbf{98^\circ}$ & $[0,98,16,114]$ & $\mathbf{3.228}$ & $\mathbf{2.1963}$\\
$114^\circ$ & $[0,114,0,114]$ & $3.639$ & $2.4761$\\
\bottomrule
\end{tabular}
\end{center}

\noindent For completeness, what abandoning the criterion on the $x$ axis would
buy. Parametrising the deviation by $\psi=\varphi_1-(\varphi_0+\varphi_2)/2$
--- the curvature of the phase profile, zero exactly when the criterion holds
--- Eq.~\eqref{eq:n3env} generalises to
$\max_t|A|=\sqrt{4\varrho^2+4\varrho|\cos\psi|+1}$, minimal at
$\psi=\pm90^\circ$. Re-optimising $\varphi_y$ and $t_0$ at each $\psi$:

\begin{center}\small
\begin{tabular}{lcccc}
\toprule
$\psi$ & $\crest_x$ & $(\crest_x\crest_y)^2$ & rel. $P_\mathrm{avg}$ & $U(\theta)$\\
\midrule
$0^\circ$ (centred) & $2.4494$ & $47.9$ & $1.00$ & $\mathbf{0.94\,\%}$\\
$45^\circ$ & $2.2836$ & $40.7$ & $1.18$ & $1.27\,\%$\\
$90^\circ$ & $\mathbf{1.8220}$ & $25.9$ & $\mathbf{1.85}$ & $3.88\,\%$\\
\bottomrule
\end{tabular}
\end{center}

\noindent Interestingly the \emph{centroid} stays at zero in all three rows:
$t_0$ and the $y$ phases can still null the odd moment of the
\emph{time-integrated} pulse area. What is lost is not the position but the
quiescence --- at $\psi=0$ the pattern stands still over the whole pulse, at
$\psi\neq0$ it breathes and only the balance over the pulse cancels. The
surviving \emph{even} part is the spread $U(\theta)$, which grows by a factor
four. For a $\pi$ pulse, $1-F\approx(\pi U/2)^2$, i.e.\ $2.2\cdot10^{-4}$
versus $3.7\cdot10^{-3}$.

\section{Mode integrals, the overlap, and the pulse area}\label{sec:area}

\subsection{The power of one spot}\label{sec:P1}
With $g_s=A_s\,u(\mathbf r-\mathbf c_s)$, a single spot carries
$\int g_s^2\dd^2r=a_sP_1$, $P_1=\int u^2\dd^2r$. For the Gaussian, substituting
$x=2r^2/w_0^2$,
\begin{equation}
P_1^\mathrm{G}=\int_0^\infty e^{-2r^2/w_0^2}2\pi r\dd r
=\frac{\pi w_0^2}{2}\int_0^\infty e^{-x}\dd x=\frac{\pi w_0^2}{2} .
\end{equation}
For the Airy mode, with $x=kr$ and
$\int_0^\infty J_\nu^2(x)\dd x/x=1/(2\nu)$ \citep{dlmf},
\begin{equation}
P_1^\mathrm{A}=\int_0^\infty\!\left(\frac{2J_1(kr)}{kr}\right)^{\!2}\!2\pi r\dd r
=\frac{8\pi}{k^2}\cdot\frac12=\frac{4\pi}{k^2},
\qquad k=\frac{3.8317}{1.4830\,w_0} .
\end{equation}
The Airy spot carries $1.20$ times the power of a Gaussian of the same
$1/e^2$ width, because of its rings.

\subsection{The overlap integral}\label{sec:overlap}
Only degenerate pairs survive the time average, and each contributes
$2A_sA_{s'}\cos(\phi_s-\phi_{s'})\,O(\varrho)$ with
$O(\varrho)=\int u(\mathbf r-\mathbf c_s)u(\mathbf r-\mathbf c_{s'})\dd^2r$ and
$\varrho=|\mathbf c_s-\mathbf c_{s'}|$.

\paragraph{Gaussian.} Completing the square with the midpoint
$\mathbf m=\tfrac12(\mathbf c_s+\mathbf c_{s'})$,
\begin{equation}
e^{-\frac{(\mathbf r-\mathbf c_s)^2}{w_0^2}}
e^{-\frac{(\mathbf r-\mathbf c_{s'})^2}{w_0^2}}
=e^{-\frac{\varrho^2}{2w_0^2}}e^{-\frac{2(\mathbf r-\mathbf m)^2}{w_0^2}}
\ \Longrightarrow\
O^\mathrm{G}=P_1^\mathrm{G}e^{-\varrho^2/2w_0^2} .
\end{equation}

\paragraph{Airy: three ingredients.}
\emph{(i) The mode is band limited.} The focal field is the Fourier transform
of the pupil field, and a uniformly illuminated circular pupil is a disc
\citep[Ch.~4,~6]{goodman2017}, \citep[Sec.~8.5]{bornwolf1999}. Hence
\begin{equation}
\tilde u(q)=\int u(|\mathbf r|)e^{-i\mathbf q\cdot\mathbf r}\dd^2r
=\frac{4\pi}{k^2}\Theta(k-q),
\end{equation}
a \emph{flat} spectrum inside the cutoff. The prefactor follows at
$\mathbf q=0$: $\tilde u(0)=2\pi\int_0^\infty u(r)r\dd r
=\frac{4\pi}{k}\int_0^\infty J_1(kr)\dd r=\frac{4\pi}{k^2}$, using
$\int_0^\infty J_1(x)\dd x=1$ \citep{dlmf}.

\emph{(ii) The integral over a disc.} The angular integral gives
$\int_0^{2\pi}e^{ix\cos\varphi}\dd\varphi=2\pi J_0(x)$ and the radial one closes
with $\int_0^XJ_0(t)t\dd t=XJ_1(X)$ (one line from
$\tfrac{\dd}{\dd z}[zJ_1(z)]=zJ_0(z)$):
\begin{equation}
\int_{|\mathbf q|<k}e^{i\mathbf q\cdot\mathbf d}\dd^2q
=2\pi\int_0^kJ_0(q\varrho)\,q\dd q=\frac{2\pi k J_1(k\varrho)}{\varrho} .
\label{eq:disc}
\end{equation}

\emph{(iii) The correlation theorem, written out.} With the convention
$\tilde u(\mathbf q)=\int u\,e^{-i\mathbf q\mathbf r}\dd^2r$ and the whole
$1/(2\pi)^2$ in the back transform \citep[Ch.~2]{goodman2017}, expand
\emph{only} the shifted factor, so that the shift becomes a phase:
\begin{align}
O(\mathbf d)&=\int\!\dd^2r\,u(\mathbf r)\,\frac{1}{(2\pi)^2}
\int\!\dd^2q\,\tilde u(\mathbf q)e^{i\mathbf q\mathbf r}e^{-i\mathbf q\mathbf d}\\
&=\frac{1}{(2\pi)^2}\int\!\dd^2q\,\tilde u(\mathbf q)e^{-i\mathbf q\mathbf d}
\underbrace{\int\!\dd^2r\,u(\mathbf r)e^{+i\mathbf q\mathbf r}}_{=\tilde u(-\mathbf q)} .
\end{align}
Interchanging the integrations is legitimate ($u$ bounded and square
integrable, $\tilde u$ of compact support, so the double integral converges
absolutely). Now the symmetry: $u$ real gives
$\tilde u(-\mathbf q)=\tilde u^{*}(\mathbf q)$, $u$ even gives
$\tilde u(-\mathbf q)=\tilde u(\mathbf q)$; both together give
$\tilde u(\mathbf q)\tilde u(-\mathbf q)=|\tilde u(\mathbf q)|^2$ and
\begin{equation}
O(\mathbf d)=\frac{1}{(2\pi)^2}\int|\tilde u(\mathbf q)|^2
e^{-i\mathbf q\cdot\mathbf d}\dd^2q .
\label{eq:corr}
\end{equation}
The sign in the exponent is immaterial because $|\tilde u|^2$ is even. At
$\mathbf d=0$, Eq.~\eqref{eq:corr} is Parseval's theorem and returns $P_1$ ---
the cheapest check of the prefactor. (``Correlation equals convolution'' is the
statement that $v(\mathbf d-\mathbf r)=v(\mathbf r-\mathbf d)$ for even $v$.)

\paragraph{Putting them together.} An indicator function is idempotent,
$\Theta^2=\Theta$, so squaring the spectrum changes nothing:
\begin{equation}
O^\mathrm{A}(\varrho)
=\frac{1}{(2\pi)^2}\left(\frac{4\pi}{k^2}\right)^{\!2}
\frac{2\pi kJ_1(k\varrho)}{\varrho}
=\frac{8\pi J_1(k\varrho)}{k^3\varrho}
=\underbrace{\frac{4\pi}{k^2}}_{P_1^\mathrm{A}}
\underbrace{\frac{2J_1(k\varrho)}{k\varrho}}_{u(\varrho)} .
\end{equation}
\textbf{The autocorrelation of the Airy mode is the mode itself}, including its
sign, which is negative between the first and the second ring. No such
statement holds for the Gaussian, whose squared spectrum is narrower by
$\sqrt2$ --- which is exactly why $O^\mathrm{G}$ carries $2w_0^2$ in the
exponent. Read in the pupil: the two spots are plane waves tilted against each
other, i.e.\ a linear phase across the pupil, and averaging it over the disc is
Eq.~\eqref{eq:disc}; the further apart, the more fringes fit and the more they
cancel. Numerically ($2048^2$ grid, \SI{40}{\micro\meter}):
$O(0)=P_1$ exactly, $P_1^\mathrm{A}=4\pi/k^2$ to \SI{1.2}{\percent}, and
$O(\varrho)=P_1^\mathrm{A}u(\varrho)$ to \SI{1.1}{\percent} of the peak over
$\varrho=0.2\dots\SI{6}{\micro\meter}$, the residual being the grid truncating
the outer rings. At $\varrho=\SI{3.305}{\micro\meter}$, $u=+0.06308$.

\subsection{Where the sinc comes from}\label{sec:sinc}
Each Fourier order of $\tilde I=\sum_dD_de^{i2\pi df_0t}$ contributes
\begin{equation}
\int_{t_0}^{t_0+T_p}\!e^{i2\pi df_0t}\dd t
=\frac{e^{i2\pi df_0(t_0+T_p)}-e^{i2\pi df_0t_0}}{i2\pi df_0} .
\end{equation}
Pulling out the phase of the window \emph{centre} makes the bracket symmetric,
and $e^{ix}-e^{-ix}=2i\sin x$ gives
\begin{equation}
=e^{i2\pi df_0(t_0+T_p/2)}\frac{\sin(\pi df_0T_p)}{\pi df_0}
=T_p\,e^{i2\pi df_0(t_0+T_p/2)}\operatorname{sinc}(df_0T_p) .
\end{equation}
The three factors read directly: $T_p$ is the window length, the exponential
says only the centre of the window matters for the phase, and the sinc is the
Fourier transform of the window itself --- integrating over a box is a
convolution with a box, and convolution in time is multiplication in frequency.
The $d=0$ term survives always, $\operatorname{sinc}(0)=1$.

\subsection{Atom-weighted averages}
$\theta(\mathbf r)$ is a field; the atom sees it through
$W(\mathbf r)=(2\pi\sigma^2)^{-1}\exp[-(\mathbf r-\mathbf r_0)^2/2\sigma^2]$
with $\sigma^2=\frac{\hbar}{2m\omega_r}\coth\frac{\hbar\omega_r}{2k_BT}$
\citep{tuchendler2008}. On a uniform grid the integral is the weighted sum
$\sum_iw_i\theta_i$, $w_i=W_i/\sum_iW_i$, since the cell area cancels; the code
uses a dedicated fine grid of $\pm4\sigma$, the global grid being far too coarse
for $\sigma\approx\SI{0.1}{\micro\meter}$. Because $\theta$ is \emph{linear} in
the $D_d$, the spatial average commutes with the sum: with
$c_d=\int WD_d\dd^2r$ computed once, a full trigger scan is one
matrix--vector product. The spread needs the map itself,
$\sigma_\theta^2=\int W\theta^2\dd^2r-\langle\theta\rangle_W^2$, and the
excitation must be averaged \emph{after} the non-linear step,
$\langle\sin^2(\theta/2)\rangle_W\neq\sin^2(\langle\theta\rangle_W/2)$.

\section{Diffraction}\label{sec:diff}

\subsection{Where the \boldmath$\sin^2$ and the $\sqrt P$ come from}
The RF drives a piezo transducer; the strain amplitude obeys
$P_\mathrm{ac}\propto S^2$, and the photoelastic effect converts strain into
index modulation \citep[Ch.~20]{salehteich2019},
\begin{equation}
\Delta n=-\tfrac12n_0^3pS\propto\sqrt{P_\mathrm{ac}} .
\label{eq:photoelast}
\end{equation}
\textbf{This is the origin of the square root} --- it is in the conversion from
power to amplitude, not in the diffraction. In the Bragg regime two waves
exchange amplitude periodically, giving $\eta=\sin^2\nu$ with
$\nu=\pi\Delta nL/(\lambda\cos\theta_\mathrm B)$, hence
\begin{equation}
\eta(P)=\alpha\sin^2\!\bigl(\kappa\sqrt P\bigr)
=\alpha\sin^2\!\left(\frac\pi2\sqrt{\frac{P}{P_\mathrm{s}}}\right),
\qquad
\kappa=\frac\pi\lambda\sqrt{\frac{M_2L}{2H}}\sqrt\zeta ,
\label{eq:etasingle}
\end{equation}
which is the standard result \citep{korpel1981,salehteich2019} and is quoted in
exactly this form by the manufacturer \citep{aaoptonotes}. Neither $M_2$ nor
$L/H$ nor the conversion $\zeta$ is known individually to useful accuracy,
which is why $\alpha$ and $P_\mathrm{s}$ are measured
\citep{mittenbuehler2024}.

\subsection{Many tones: coupled waves and the exact reduction}\label{sec:coupled}
Write the field as one undiffracted wave plus $N$ diffracted ones,
$E=\Au e^{i(\mathbf k_\mathrm u\mathbf r-\omega_0t)}
+\sum_{n}A_ne^{i(\mathbf k_n\mathbf r-\omega_nt)}$, with
$\omega_n=\omega_0+\Omega_n$. \emph{$\Au$ carries no number --- it belongs to
no tone; $A_n$ is the beam diffracted by tone $n$, and every $\sum_n$ below
runs over tones and never includes $\Au$.} Substituting into the Helmholtz
equation with $n^2\approx n_0^2+2n_0\sum_n\Delta n_n\cos(\mathbf K_n\mathbf r
-\Omega_nt+\varphi_n)$ and collecting terms at $\omega_n$ gives
\citep{kogelnik1969,laude2003}
\begin{align}
\frac{\dd A_n}{\dd z}&=-i\kappa_ne^{i\varphi_n}\Au\,e^{i\Delta k_nz},
&\frac{\dd\Au}{\dd z}&=-i\sum_n\kappa_ne^{-i\varphi_n}A_n\,e^{-i\Delta k_nz},
\label{eq:coupledgen}
\end{align}
with $\kappa_n=\kappa\sqrt{P_n}$ and
$\Delta k_n=(\mathbf k_\mathrm u+\mathbf K_n-\mathbf k_n)\cdot\hat{\mathbf z}$.
Three readings:

\begin{enumerate}\itemsep1pt
\item \textbf{The tone phases drop out.} $\tilde A_n=A_ne^{-i\varphi_n}$
      removes $\varphi_n$ from both equations, and $|A_n|=|\tilde A_n|$: the
      diffraction efficiency is \emph{exactly} phase independent, and
      $\varphi_n$ reappears only as the phase of the diffracted wave, i.e.\ in
      the interference pattern. This is why the crest factor --- a property of
      the \emph{sum}, which only the broadband amplifier sees --- cannot appear
      in the diffraction at all.
\item \textbf{Phase matching is assumed here, not later.} Setting
      $\Delta k_n=0$ leaves the working equations; $z$ is measured in units of
      the interaction length, so the exit is $z=1$ and $\kappa_n$ is the
      accumulated coupling.
\item \textbf{Unequal $\Delta k_n$ break the reduction below}, because carrying
      $e^{-i\Delta k_nz}$ through the derivative creates a new combination with
      weights $\kappa_n\Delta k_n$, and the hierarchy closes only if all
      $\Delta k_n$ are equal (then one gets the detuned two-wave solution) or
      zero.
\end{enumerate}

\paragraph{Bright mode.} With $\kappa_\mathrm{eff}=\sqrt{\sum_n\kappa_n^2}$ and
$B=\kappa_\mathrm{eff}^{-1}\sum_n\kappa_nA_n$ one gets
$\dd\Au/\dd z=-i\kappa_\mathrm{eff}B$ and
$\dd B/\dd z=-i\kappa_\mathrm{eff}\Au$: a closed two-wave system.

\paragraph{Dark modes stay dark.} Any $C=\sum_nc_nA_n$ with
$\sum_nc_n\kappa_n=0$ obeys $\dd C/\dd z=0$ and $C(0)=0$, so $C\equiv0$. Of the
$N$ diffracted degrees of freedom, $N-1$ are never excited --- \emph{this} is
what makes the reduction exact rather than approximate.

\paragraph{Solve and return.} $\Au=\cos(\kappa_\mathrm{eff}z)$,
$B=-i\sin(\kappa_\mathrm{eff}z)$. Integrating $\dd A_n/\dd z=-i\kappa_n\Au$ with
$A_n(0)=0$, the \emph{same} integral appears for every $n$:
\begin{equation}
A_n(z)=-i\kappa_n\!\int_0^z\!\Au(z')\dd z'
=-i\frac{\kappa_n}{\kappa_\mathrm{eff}}\sin(\kappa_\mathrm{eff}z) ,
\end{equation}
so the orders differ only by the prefactor. Summing at $z=1$,
$\eta_\mathrm{tot}=\sum_n|A_n|^2=\sin^2\kappa_\mathrm{eff}=1-|\Au|^2$, and ---
the step where the sum moves under the root ---
\begin{equation}
\kappa_\mathrm{eff}=\sqrt{\textstyle\sum_n\kappa_n^2}
=\sqrt{\textstyle\sum_n\kappa^2P_n}=\kappa\sqrt{P_\mathrm{tot}} .
\end{equation}
The squares of the couplings \emph{are} the powers, so their quadratic sum is
the ordinary sum of powers under one root. Hence
\begin{empheq}[box=\fbox]{equation}
\eta_\mathrm{tot}=\alpha\sin^2\!\bigl(\kappa\sqrt{P_\mathrm{tot}}\bigr),
\qquad
\eta_n=\frac{P_n}{P_\mathrm{tot}}\eta_\mathrm{tot} .
\label{eq:multitone}
\end{empheq}
Energy is conserved exactly, and for $N=1$ this reduces to
\eqref{eq:etasingle}. The same result, derived independently and verified
experimentally on the same class of device, is reported in
\citep{mittenbuehler2024}; a treatment with optimised tone phases is
\citep{antonov2007}.

\subsection{Why the naive sum is wrong here}
$\sum_n\sin^2(\kappa\sqrt{P_n})$ is not wrong physics but the small-signal
limit taken in the wrong place: for $\kappa\sqrt P\ll1$ both expressions tend
to $\kappa^2P_\mathrm{tot}$. They separate exactly where we work:

\begin{center}\small
\begin{tabular}{cccc}
\toprule
$P_\mathrm{tot}$ & $\kappa\sqrt{P_\mathrm{tot}}$ & $\eta_\mathrm{tot}$ exact
 & naive sum (3 tones)\\
\midrule
\SI{5}{\milli\watt} & $0.066$ & $0.00432$ & $0.00433$ \ ($+0.10\,\%$)\\
\SI{50}{\milli\watt} & $0.208$ & $0.0427$ & $0.0431$ \ ($+0.97\,\%$)\\
\textbf{\SI{0.664}{\watt}} & $0.758$ & $\mathbf{0.473}$ & $0.539$ \ ($\mathbf{+14.0\,\%}$)\\
\bottomrule
\end{tabular}
\end{center}

\subsection{Validity}
Equation~\eqref{eq:multitone} is exact at any coupling strength (checked
against direct integration to $10^{-13}$, including beyond the first maximum
where back-conversion sets in). What is assumed is:

\begin{enumerate}\itemsep1pt
\item \textbf{Common phase matching.} For $\Delta k\neq0$ the two-wave solution
      becomes
      $\eta(\xi)=\nu^2(\nu^2+\xi^2)^{-1}\sin^2\sqrt{\nu^2+\xi^2}$ with
      $\nu=\kappa\sqrt P$, $\xi=\Delta kL/2$ \citep{kogelnik1969}. Calibrating
      $\xi$ against the specified \SI{36}{\mega\hertz} bandwidth (half
      efficiency at the band edge, $\xi_{1/2}=1.3625$ at $\nu=0.758$) gives
      $\xi=1.3625\,\delta f/\SI{18}{\mega\hertz}$ and hence
      \SI{0.007}{\percent} loss for the \SI{0.37}{\mega\hertz} span of the
      working point, \SI{1.8}{\percent} for \SI{6}{\mega\hertz}, and
      \SI{15}{\percent} for \SI{18}{\mega\hertz}.
\item \textbf{Bragg regime, two waves per grating}; higher orders are strongly
      mismatched.
\item \textbf{Linear transducer.} The instantaneous strain carries the crest
      factor as well, so a high $\crest$ also brings on acoustic nonlinearity
      --- a second channel, independent of the amplifier, pointing the same
      way, not quantified here.
\item \textbf{Isotropic diffraction.} The device diffracts anisotropically in a
      birefringent crystal; the isotropic result agrees qualitatively
      \citep{chang1995}.
\item \textbf{Independent axes.} Measurement shows a cross term,
      $\eta_{ij}=\eta^{(x)}_i\eta^{(y)}_j\Xi_{ij}$ with
      $\Xi=(1-\xi_0(\nu^{(x)}_i-\xi_x)^2(\nu^{(y)}_j-\xi_y)^2)^2$, falling to
      $0.80$ in the corners of the scan range \citep{mittenbuehler2024}. Over
      the \SI{0.37}{\mega\hertz} span used here $\Xi$ is common to all spots.
\end{enumerate}

\subsection{The crest factor in the diffraction}
Combining the back-off rule with \eqref{eq:multitone},
\begin{equation}
\eta_\mathrm{axis}(\crest)=\alpha\sin^2\!\left(\frac{\pi}{2\crest}
\sqrt{\frac{P_\mathrm{1dB}}{P_\mathrm{s}}}\right) .
\end{equation}
With the test-sheet parameters ($\alpha=1$, $P_\mathrm{s}=\SI{2.85}{\watt}$,
$P_\mathrm{1dB}=\SI{3.98}{\watt}$) this is $\sin^2(1.8559/\crest)$:

\begin{center}\small
\begin{tabular}{lccc}
\toprule
& $\crest$ & $P_\mathrm{avg}$ & $\eta_\mathrm{axis}$\\
\midrule
theoretical minimum & $\sqrt2=1.414$ & \SI{1.99}{\watt} & $0.935$\\
$x$ with $\psi=90^\circ$ (criterion abandoned) & $1.822$ & \SI{1.20}{\watt} & $0.725$\\
$y$, $b=98^\circ$ & $2.196$ & \SI{0.83}{\watt} & $0.559$\\
$\mathbf x$, \textbf{centred}, $\mathbf{N=3}$ & $\mathbf{2.449}$ & \textbf{\SI{0.66}{\watt}} & $\mathbf{0.472}$\\
$y$, ramp & $2.827$ & \SI{0.50}{\watt} & $0.373$\\
\bottomrule
\end{tabular}
\end{center}

\noindent Had both axes reached $\crest=\sqrt2$, the profile efficiency would
be \SI{87}{\percent} instead of \SI{26.4}{\percent} --- a factor $3.3$. That
factor is the price of the crest factor, and Sec.~\ref{sec:n3} says how much of
it is negotiable: on the $x$ axis none, as long as three tones are used and the
pattern must stay centred.

\bibliographystyle{unsrtnat}
\bibliography{refs}
\end{document}
